\documentclass[12pt,a4paper]{article}

\usepackage[utf8]{inputenc}
\usepackage[T2A]{fontenc}
\usepackage[russian]{babel}
\usepackage{amsmath,amssymb,amsthm}
\usepackage{geometry}
\usepackage{hyperref}
\geometry{margin=2.2cm}
\hypersetup{colorlinks=true,linkcolor=blue,urlcolor=blue}

\title{Алгоритм KA008: робастное равномерное самораспределение\\
и построение виртуальной антенной решётки в 3D (формулы)}
\author{Конспект по статье: \textit{Robust Self-Distribution of Homogeneous Robotic Teams...} (Boyko et al.)\\
Файл-источник: \texttt{BooksDocuments/KA008.pdf}.}
\date{}

\begin{document}
\maketitle

\begin{abstract}
Этот документ кратко и последовательно выписывает алгоритм из \texttt{KA008.pdf} со всеми ключевыми формулами:
определение видимости/информантов, ``зазоры'' $g_i^\pm$, функция $\Xi(\cdot)$, скользящая переменная $\Sigma_i$,
переключающее (sliding-mode) управление $u_i(t)=u_i^{\langle-\mathrm{sgn}\,\Sigma_i(t)\rangle}$,
а также модификации для (i) стационарного равномерного распределения с шагом $w$, (ii) пейсмейкера, (iii) динамического якоря
без начальной информационной связности. В конце приведена композиция по осям для построения 3D ``вертикальной антенны'' над маяком
из раздела VI (computer simulations).
\end{abstract}

\section{Постановка (скалярное состояние)}

Есть $N\ge 2$ одинаковых роботов. Для робота $i$ вводится скалярное состояние $y_i(t)\in\mathbb{R}$ и скорость
$v_i(t):=\dot y_i(t)$. Роботы могут иметь сложную (в т.ч. неопределённую) динамику второго порядка
\begin{equation}
  \ddot y_i = f_i\bigl[t,\,Y(t),\,V(t),\,u_i(t),\,\mu_i(t)\bigr], \qquad u_i\in\Omega_i,
\end{equation}
где $Y(t)=\{y_i(t)\}$, $V(t)=\{v_i(t)\}$, а $\mu_i$ агрегирует возмущения/неучтённую динамику (в статье это формулы (5)--(7)).
Ключевое инженерное допущение: у каждого робота есть два управляющих воздействия
$u_i^{\langle+1\rangle},u_i^{\langle-1\rangle}\in\Omega_i$, которые гарантированно дают разгон/торможение (Asm.~4 статьи):
\begin{equation}
  \pm f_i\bigl[t,Y,V,u_i^{\langle\pm1\rangle},\mu_i(t)\bigr] \ge e_{u,i} > 0.
  \label{eq:asm4}
\end{equation}

\subsection*{Цель}

Состояния должны прийти к равномерной решётке с заданным шагом $w\ge 0$ в одной из постановок:
\begin{itemize}
  \item (1) \textbf{Steady distribution}: $y_i\to y^\star + i w$ (для некоторого перечисления), скорости $\to 0$.
  \item (2) \textbf{Dynamic anchor}: один край решётки задаётся якорем $y^\dagger(t)$ (крайний ``нижний/верхний'' узел).
  \item (3) \textbf{Pacemaker}: $y^\dagger(t)$ является одним из узлов решётки (фиктивный робот $-1$).
\end{itemize}

\section{Видимость и ``зазоры'' $g_i^\pm$}

Робот видит только ближайшего соседа \emph{по направлению} в пределах радиуса видимости $R_{\mathrm{vis}}$:
вверх (если $y_j>y_i$) и вниз (если $y_j<y_i$), причём ``прямая видимость'' блокируется промежуточными роботами (Asm.~1).
Поэтому для каждого $i$ определяются верхний и нижний ``зазоры'' с параметром $\varsigma\ge 0$:
\begin{align}
g_i^+(t\mid \varsigma) &:= 
\begin{cases}
  y_j(t)-y_i(t), & \text{если есть видимый сосед $j$ выше},\\
  \varsigma, & \text{иначе},
\end{cases}
\label{eq:gplus}\\
g_i^-(t\mid \varsigma) &:=
\begin{cases}
  y_i(t)-y_j(t), & \text{если есть видимый сосед $j$ ниже},\\
  \varsigma, & \text{иначе}.
\end{cases}
\label{eq:gminus}
\end{align}

Для целевой решётки требуется $w<R_{\mathrm{vis}}$ (Asm.~3).

\section{Выбор функции $\Xi(\cdot)$}

Во всех правилах используется гладкая монотонная функция $\Xi:[0,R_{\mathrm{vis}}]\to\mathbb{R}$:
\begin{equation}
  \Xi'(d) > 0 \ \forall d\in[0,R_{\mathrm{vis}}], \qquad \Xi(0)=0.
  \label{eq:Xi}
\end{equation}
Примеры из статьи: $\Xi(d)=k d$, $\Xi(d)=k\arctan d$, $\Xi(d)=\frac{k d}{1+d}$, либо $\Xi(d)=\tanh(d)$ при подходящем масштабировании.

\section{Правило управления 1 (равномерное распределение с шагом $w$)}

Определяется скользящая (sliding) переменная:
\begin{equation}
  \Sigma_i(t) := v_i(t) - \Xi\!\bigl(g_i^+(t\mid w)\bigr) + \Xi\!\bigl(g_i^-(t\mid w)\bigr).
  \label{eq:Sigma1}
\end{equation}

Далее применяется переключающее управление (формулы (12)--(13) статьи):
\begin{equation}
  u_i(t) := u_i^{\langle -\mathrm{sgn}\,\Sigma_i(t)\rangle},
  \label{eq:u1}
\end{equation}
где $\mathrm{sgn}(\cdot)$ --- знак, а $u_i^{\langle\pm1\rangle}$ --- ``разгон/торможение'' из \eqref{eq:asm4}.
Движение рассматривается в смысле Филиппова (из-за разрывности по $\Sigma_i=0$).

\paragraph{Смысл.}
Если сверху зазор больше (через $\Xi(g_i^+)$), а снизу меньше, то $\Sigma_i$ меняет знак так, что робот
переключается между ускорением и торможением, выравнивая зазоры до $w$ при одновременном затухании скорости.

\section{Пейсмейкер (control rule 2)}

В задачах с пейсмейкером вводится фиктивный робот $-1$ со состоянием $y^\dagger(t)$.
Истинные роботы не отличают его от остальных (Asm.~2) и применяют то же правило \eqref{eq:gplus}--\eqref{eq:u1},
но вычисляя $g_i^\pm$ с учётом $y^\dagger$ как ещё одного ``пира'' (формально: множества видимых объектов расширяются).
Теорема 2 статьи даёт условия сходимости при ограниченных скорости/ускорении пейсмейкера.

\section{Динамический якорь без начальной связности (control rule 3)}

Для ``нижнего якоря'' $y^\dagger(t)$ (край решётки) вводится модификация:
\begin{equation}
  \Sigma_i(t) := v_i(t) - \Xi\!\bigl(g_i^+(t\mid w)\bigr) + \Xi\!\bigl(g_i^-(t\mid R_{\mathrm{vis}})\bigr),
  \label{eq:Sigma3}
\end{equation}
где верхний зазор дополнительно ограничивается:
\begin{equation}
  g_i^+(t\mid w) := \min\{g_i^+(t\mid R_{\mathrm{vis}}),\ w\}.
  \label{eq:gplus_clip}
\end{equation}
Тогда управление снова $u_i(t)=u_i^{\langle -\mathrm{sgn}\,\Sigma_i(t)\rangle}$.
Эта модификация позволяет \emph{создавать информационную связность ``из ничего''} (Remark 7): даже если сначала роботы далеко
по состоянию (нет рёбер видимости), они со временем сближаются с якорем и формируют решётку.

\section{Композиция по осям для 3D ``виртуальной антенны'' (раздел VI)}

В разделе VI статья рассматривает 3D задачу: над маяком (beacon) $p(t)$ с нулевой высотой построить вертикальную колонну:
горизонтально $x_i,y_i\to 0$ в общей системе $F$ с началом в $p$, а высоты выстроить с шагом $w_a$ и нижний робот на высоте $h$,
где $h\ge w_a>0$.

Делается разбиение ресурсов по осям (управления и дальности):
\begin{equation}
  \sum_{b\in\{x,y,z\}} u_b^2 \le u^2,\qquad \sum_{b\in\{x,y,z\}} R_b^2 \le R_{\mathrm{sen}}^2,\qquad R_z > w_a,\ h.
\end{equation}
Для каждой оси выбирается $\Xi_b:[0,R_b]\to[0,u_b)$, удовлетворяющая \eqref{eq:Xi} (с $R_{\mathrm{vis}}:=R_b$).
Управление по координатам формируется как вектор $u_i(t)=(u_i^x,u_i^y,u_i^z)$.

\subsection*{Горизонталь ($x,y$): консенсус/стягивание к маяку}

Для $b\in\{x,y\}$ применяется правило \eqref{eq:u1}--\eqref{eq:Sigma1} с \textbf{$w:=0$} и $\Xi:=\Xi_b$ к
относительным координатам $b_j(t)-b_i(t)$ всех видимых объектов $j$, включая маяк $p$ (если маяк видим).
То есть по сути это ``консенсус к нулю'' по каждой горизонтальной координате.

\subsection*{Вертикаль ($z$): решётка с якорем и шагом $w_a$}

Для координаты $z$ применяется модифицированное правило (в статье: ``control rule 3 with $w:=w_a$ and $R_{\mathrm{vis}}:=R_z$''),
то есть \eqref{eq:u1} и \eqref{eq:Sigma3}--\eqref{eq:gplus_clip} с $\Xi:=\Xi_z$ и шагом $w_a$.
Дополнительно в разделе VI сказано: ``координата маяка предварительно умножается на $w_a/h$'' (масштабирование якоря по высоте),
чтобы согласовать требование ``нижний робот на высоте $h$'' с шагом $w_a$.

\subsection*{Инженерная реализация (Webots/Crazyflie)}

В конце раздела VI указано, что для реалистичной модели квадрокоптера вычисляемые ``силы/ускорения'' по осям затем
преобразуются в управляющие сигналы модели через банк PD/PID регуляторов (детали опущены в статье).

\section*{Итого: алгоритм в псевдокоде (1D)}

Для каждого робота $i$ в каждый момент времени:
\begin{enumerate}
  \item измерить $y_i$, $v_i$, а также определить видимого соседа сверху и снизу;
  \item вычислить $g_i^+(t\mid \varsigma)$ и $g_i^-(t\mid \varsigma)$ по \eqref{eq:gplus}--\eqref{eq:gminus}
        (где $\varsigma=w$ для rule~1/2, а в rule~3 используется $R_{\mathrm{vis}}$ внизу и ограничение \eqref{eq:gplus_clip} сверху);
  \item вычислить $\Sigma_i$ по \eqref{eq:Sigma1} (rule~1/2) или \eqref{eq:Sigma3} (rule~3);
  \item выдать управляющее воздействие $u_i := u_i^{\langle -\mathrm{sgn}\,\Sigma_i\rangle}$.
\end{enumerate}

\end{document}

